% Part: methods
% Chapter: proofs
% Section: example-1

\documentclass[../../../include/open-logic-section]{subfiles}

\begin{document}

\olfileid[pt]{mth}{prf}{ex1}

\olsection{Um Exemplo} 

Nosso primeiro exemplo é o seguinte fato simples sobre uniões e
interseções de conjuntos.  Ele ilustrará o desdobramento de definições,
provas de conjunções, de afirmações universais, e a prova por casos.

\begin{prop}
Para quaisquer conjuntos $A$, $B$ e $C$, $A \cup (B \cap C) = (A \cup B)
\cap (A \cup C)$
\end{prop}

Vamos prová-la!{}

\begin{proof}
Queremos mostrar que, para quaisquer conjuntos $A$, $B$ e $C$, $A \cup (B \cap
C) = (A \cup B) \cap (A \cup C)$
\begin{quote}
Primeiro desdobramos a definição de ``$=$'' no enunciado da
proposição. Lembre-se de que provar que conjuntos são idênticos significa mostrar que os
conjuntos têm os mesmos !!{element}s. Isto é, todos os !!{element}s de $A \cup
(B \cap C)$ são também !!{element}s de $(A \cup B) \cap (A \cup C)$, e
vice-versa.  O ``vice-versa'' significa que também todo !!{element} de $(A
\cup B) \cap (A \cup C)$ deve ser !!a{element} de $A \cup (B \cap
C)$.  Assim, ao desdobrar a definição, vemos que temos que provar uma
conjunção.  Vamos registrar isso:
\end{quote}
Por definição, $A \cup (B \cap C) = (A \cup B) \cap (A \cup C)$ sse
todo !!{element} de $A \cup (B \cap C)$ é também !!a{element} de $(A
\cup B) \cap (A \cup C)$, e todo !!{element} de $(A \cup B) \cap (A
\cup C)$ é !!a{element} de $A \cup (B \cap C)$.
\begin{quote}
Como isso é uma conjunção, devemos provar cada conjuntiva
separadamente. Vamos começar pela primeira: vamos provar que todo
!!{element} de $A \cup (B \cap C)$ é também !!a{element} de $(A
\cup B) \cap (A \cup C)$.

Esta é uma afirmação universal, e por isso consideramos um !!{element}
arbitrário de $A \cup (B \cap C)$ e mostramos que ele também deve ser !!a{element} de
$(A \cup B) \cap (A \cup C)$. Vamos escolher uma variável para chamar esse
!!{element} arbitrário, digamos,~$z$.  Nossa prova continua:
\end{quote}
Primeiro, provamos que todo !!{element} de $A \cup (B \cap C)$ é também
!!a{element} de $(A \cup B) \cap (A \cup C)$. Seja $z \in A \cup (B
\cap C)$. Temos que mostrar que $z \in (A \cup B) \cap (A \cup C)$.
\begin{quote}  
Agora é hora de desdobrar a definição de $\cup$ e~$\cap$. Por
exemplo, a definição de $\cup$ é: $A \cup B = \Setabs{z}{z \in A
  \text{ ou } z \in B}$.  Quando aplicamos a definição a ``$A \cup (B
\cap C)$,'' o papel do ``$B$'' na definição passa a ser desempenhado por
``$B \cap C$,'' de modo que $A \cup (B \cap C) = \Setabs{z}{z \in A \text{ ou }
  z \in B \cap C}$.  Então nossa suposição de que $z \in A \cup (B \cap C)$
equivale a: $z \in \Setabs{z}{z \in A \text{ ou } z \in B \cap
  C}$. E $z \in \Setabs{z}{\dots z\dots}$ sse \dots $z$ \dots, isto é,
neste caso, $z \in A$ ou $z \in B \cap C$.
\end{quote}
Pela definição de $\cup$, ou $z \in A$ ou $z \in B \cap C$.
\begin{quote}
Como isso é uma disjunção, será útil aplicar a prova por
casos. Tomamos os dois casos, e mostramos que, em cada um, a
conclusão que buscamos (a saber, ``$z \in (A \cup B) \cap (A \cup
C)$'') se obtém.
\end{quote}
Caso 1: Suponha que $z \in A$.
\begin{quote}
Não há muito mais com que trabalhar a partir das nossas suposições. Então vamos
olhar para o que temos com que trabalhar na conclusão. Queremos mostrar
que $z \in (A \cup B) \cap (A \cup C)$. Com base na definição de
$\cap$, se queremos mostrar que $z \in (A \cup B) \cap (A \cup C)$,
temos que mostrar que ele está tanto em $(A \cup B)$ quanto em $(A \cup C)$. Mas $z
\in A \cup B$ sse $z \in A$ ou $z \in B$, e já temos (como a
suposição do caso~1) que $z \in A$. Pelo mesmo
raciocínio---trocando $C$ por $B$---$z \in A \cup C$. Esse argumento
foi na direção inversa, então vamos registrar nosso raciocínio na
direção necessária à nossa prova.
\end{quote}
Como $z \in A$, $z \in A$ ou $z \in B$, e logo, pela definição
de~$\cup$, $z \in A \cup B$. De modo semelhante, $z \in A \cup C$.  Mas isso
significa que $z \in (A \cup B) \cap (A \cup C)$, pela definição
de~$\cap$.
\begin{quote}
Isso completa o primeiro caso da prova por casos. Agora queremos
derivar a conclusão no segundo caso, em que $z \in B \cap C$.
\end{quote}
Caso 2: Suponha que $z \in B \cap C$.
\begin{quote}
Novamente, estamos trabalhando com a interseção de dois conjuntos. Vamos aplicar
a definição de~$\cap$:
\end{quote}
Como $z \in B \cap C$, $z$ deve ser !!a{element} tanto de $B$ quanto de $C$, pela
definição de~$\cap$.
\begin{quote}
É hora de olhar de novo para a nossa conclusão. Temos que mostrar que $z$ está
tanto em $(A \cup B)$ quanto em $(A \cup C)$. E, novamente, a solução é
imediata.
\end{quote}
Como $z \in B$, $z \in (A \cup B)$. Como $z \in C$, também $z \in (A
\cup C)$.  Então, $z \in (A \cup B) \cap (A \cup C)$.
\begin{quote}
Aqui aplicamos as definições de $\cup$ e $\cap$ novamente, mas como
já relembramos essas definições, e já mostramos que, se
$z$ está em um de dois conjuntos, ele está na união deles, não precisamos ser tão
explícitos no que fizemos.

Completamos o segundo caso da prova por casos, então agora podemos
afirmar nossa primeira conclusão.
\end{quote}
Então, se $z \in A \cup (B \cap C)$, então $z \in (A \cup B) \cap (A \cup C)$.
\begin{quote}
Agora queremos apenas mostrar a outra direção, que todo !!{element}
de $(A \cup B) \cap (A \cup C)$ é !!a{element} de $A \cup (B \cap
C)$. Como antes, provamos essa afirmação universal supondo que temos um
!!{element} arbitrário do primeiro conjunto e mostrando que ele deve estar no
segundo conjunto. Vamos declarar o que estamos prestes a fazer.
\end{quote}
Agora, suponha que $z \in (A \cup B) \cap (A \cup C)$. Queremos mostrar
que $z \in A \cup (B \cap C)$.
\begin{quote}
Estamos agora trabalhando a partir da hipótese de que $z \in (A \cup B) \cap (A
\cup C)$. Esperamos que não seja confuso demais estarmos usando o
mesmo~$z$ aqui que na primeira parte da prova.  Quando terminamos
aquela parte, todas as suposições que fizemos ali deixaram de estar em
vigor, então agora podemos fazer novas suposições sobre o que $z$ é.  Se isso
for confuso para você, basta substituir $z$ por uma variável diferente no
que segue.

Sabemos que $z$ está tanto em $A \cup B$ quanto em $A \cup C$, pela definição
de~$\cap$. E, pela definição de $\cup$, podemos desdobrar isso ainda mais
em: ou $z \in A$ ou $z \in B$, e também ou $z \in A$ ou $z \in
C$. Isso parece uma prova por casos de novo---exceto que o ``e'' a torna
confusa. Você poderia pensar que isso equivale a haver três
possibilidades: $z$ está ou em $A$, $B$ ou $C$. Mas isso seria um
erro.  Temos que ter cuidado, então vamos considerar cada disjunção por
vez.
\end{quote}
Pela definição de $\cap$, $z \in A \cup B$ e $z \in A \cup C$. Pela
definição de $\cup$, $z \in A$ ou $z \in B$. Distinguimos casos.
\begin{quote}
Como estamos nos concentrando na primeira disjunção, ainda não obtivemos nossa
segunda disjunção (do desdobramento de $A \cup C$). De fato, ainda não
precisamos dela. O primeiro caso é $z \in A$, e !!a{element} de um conjunto é
também !!a{element} da união desse conjunto com qualquer outro. Então o caso~1
é fácil:
\end{quote}
Caso 1: Suponha que $z \in A$. Segue-se que $z \in A \cup (B \cap
C)$.
\begin{quote}
Agora, para o segundo caso, $z \in B$. Aqui desdobraremos o segundo
$\cup$ e faremos outra prova por casos:
\end{quote}
Caso 2: Suponha que $z \in B$. Como $z \in A \cup C$, ou $z \in
A$ ou $z \in C$. Distinguimos os casos ainda mais:

Caso 2a: $z \in A$. Então, novamente, $z \in A \cup (B \cap C)$.
\begin{quote}
Ok, isso foi um pouco estranho. Na verdade, não precisamos da suposição
de que~$z \in B$ para este caso, mas tudo bem.
\end{quote}
Caso 2b: $z \in C$. Então $z \in B$ e $z \in C$, de modo que $z \in B \cap C$,
e consequentemente, $z \in A \cup (B \cap C)$.
\begin{quote}
Isso conclui ambas as provas por casos, e assim terminamos a segunda
metade.
\end{quote}
Então, se $z \in (A \cup B) \cap (A \cup C)$, então $z \in A \cup (B \cap C)$.
\end{proof}

\end{document}
